\chapter{Concluzii}

\section{Rezultate ob\c tinute}

La finalul acestui proiect pot spune c\u a obiectivul principal a fost
atins: Gym-Connect este o aplica\c tie web func\c tional\u a, care
rezolv\u a exact problema de la care am pornit -- gestionarea ordonat\u a
a antrenamentelor de for\c t\u a \c si urm\u arirea progresului pe termen
lung, totul gratuit \c si accesibil direct din browser.

Concret, aplica\c tia implementeaz\u a toate func\c tionalit\u a\c tile
propuse \^ in capitolul introductiv. Utilizatorul \^ i\c si poate crea
cont \c si autentifica \^ in siguran\c t\u a prin Clerk, \^ i\c si poate
defini rutine de antrenament personalizate cu exerci\c tii, seturi \c si
intervale de repeti\c tii \c tint\u a, \c si poate \^ inregistra sesiuni
de antrenament \^ in timp real. Modulul de antrenament ofer\u a un timer
de odihn\u a \^ intre seturi, un calculator de distribu\c tie a pl\u acilor
pe bar\u a \c si suport pentru indicatorii de intensitate RIR \c si 1RM
estimat. \^ In plus, toate datele introduse de un utilizator sunt
izolate complet de ale celorlal\c ti, prin legarea fiec\u arei
\^ inregistr\u ari de identificatorul Clerk.

Cea mai valoroas\u a parte a aplica\c tiei r\u am\^ ane, \^ in opinia mea,
modulul de analitic\u a. Acesta nu doar afi\c seaz\u a un istoric brut
al antrenamentelor, ci transform\u a datele \^ in informa\c tie util\u a:
calculeaz\u a volumul s\u apt\u am\^ anal de antrenament, estimeaz\u a
evolu\c tia for\c tei prin 1RM \c si, folosind regresia liniar\u a,
proiecteaz\u a o tendin\c t\u a pe urm\u atoarele s\u apt\u am\^ ani.
Func\c tia de detec\c tie automat\u a a platourilor avertizeaz\u a
utilizatorul atunci c\^ and progresul stagneaz\u a, suger\^ and momentul
potrivit pentru un deload. Tocmai aceste func\c tionalit\u a\c ti --
absente din versiunile gratuite ale aplica\c tiilor concurente
analizate \^ in Capitolul 2 -- reprezint\u a contribu\c tia original\u a
principal\u a a acestei lucr\u ari.

Din punct de vedere tehnic, proiectul a confirmat c\u a o stiv\u a
modern\u a, omogen\u a TypeScript at\^ at pe frontend c\^ at \c si pe
backend, permite dezvoltarea rapid\u a a unei aplica\c tii complete,
cu un cod u\c sor de \^ intre\c tinut \c si extins.

\section{Limit\u arile solu\c tiei actuale}

De\c si aplica\c tia \^ i\c si \^ indepline\c ste scopul, sunt con\c stient
de o serie de limit\u ari pe care le-am identificat pe parcurs.

\^ In primul r\^ and, predic\c tia evolu\c tiei 1RM se bazeaz\u a pe
regresie liniar\u a, un model simplu care presupune o cre\c stere
constant\u a \^ in timp. \^ In realitate, progresul \^ in for\c t\u a
nu este liniar -- la \^ inceput cre\c sterile sunt rapide, apoi
\^ incetinesc. Un model mai sofisticat ar oferi predic\c tii mai
realiste, mai ales pe orizonturi lungi.

\^ In al doilea r\^ and, biblioteca de exerci\c tii este, \^ in acest
moment, predefinit\u a \^ in aplica\c tie. Utilizatorul nu poate ad\u auga
exerci\c tii proprii care nu se reg\u asesc \^ in list\u a, ceea ce
poate fi o limitare pentru sportivii cu rutine mai pu\c tin
conven\c tionale.

\^ In al treilea r\^ and, aplica\c tia func\c tioneaz\u a exclusiv online.
Nu exist\u a un mod offline -- dac\u a utilizatorul pierde conexiunea
la internet \^ in timpul antrenamentului, nu \^ i\c si poate salva
sesiunea. Pentru o aplica\c tie folosit\u a \^ intr-o sal\u a de for\c t\u a,
unde semnalul poate fi slab, acesta este un dezavantaj real.

\^ In fine, dependen\c ta de servicii externe -- Clerk pentru
autentificare \c si Neon pentru baza de date -- \^ inseamn\u a c\u a
disponibilitatea aplica\c tiei depinde de disponibilitatea acestor
servicii. \^ Intr-un context de produc\c tie serioas\u a, aceast\u a
dependen\c t\u a ar trebui evaluat\u a cu aten\c tie.

\section{Direc\c tii de dezvoltare ulterioar\u a}

Pornind de la limit\u arile de mai sus, exist\u a mai multe direc\c tii
\^ in care Gym-Connect ar putea evolua.

O prim\u a \^ imbun\u at\u a\c tire natural\u a ar fi \^ inlocuirea
regresiei liniare cu un model de predic\c tie mai apropiat de
realitatea fiziologic\u a a antrenamentului, eventual o curb\u a
logaritmic\u a sau un model care \c tine cont de istoricul individual
al fiec\u arui utilizator.

O a doua direc\c tie ar fi ad\u augarea posibilit\u a\c tii de a
defini exerci\c tii personalizate \c si chiar categorii de exerci\c tii
proprii, oferind astfel mai mult\u a flexibilitate.

Transformarea aplica\c tiei \^ intr-o \textbf{Progressive Web App}
(PWA) ar rezolva problema modului offline: utilizatorul \c si-ar
putea \^ inregistra antrenamentul local, urm\^ and ca datele s\u a se
sincronizeze automat la reconectare. Aceasta ar fi, probabil,
cea mai util\u a \^ imbun\u at\u a\c tire din perspectiva utiliz\u arii
reale \^ in sal\u a.

Pe termen mai lung, mi-ar pl\u acea s\u a explorez integrarea unor
func\c tionalit\u a\c ti sociale -- posibilitatea de a urm\u ari
progresul prietenilor sau de a compara rezultate -- \c si chiar
sugestii inteligente de antrenament, generate pe baza istoricului
\c si a obiectivelor utilizatorului.

Indiferent de direc\c tia aleas\u a, consider c\u a fundamentul
construit \^ in cadrul acestei lucr\u ari -- o arhitectur\u a curat\u a,
modular\u a \c si bine documentat\u a -- ofer\u a o baz\u a solid\u a
pentru oricare dintre aceste dezvolt\u ari viitoare.
